\documentclass[UTF8,11pt,fontset=none]{ctexart}
\usepackage[a4paper,margin=23mm,headheight=15pt,footskip=11mm]{geometry}
\usepackage{fontspec}
\IfFontExistsTF{TeX Gyre Termes}{\setmainfont{TeX Gyre Termes}}{\setmainfont{lmroman10-regular.otf}[BoldFont=lmroman10-bold.otf,ItalicFont=lmroman10-italic.otf,BoldItalicFont=lmroman10-bolditalic.otf]}
\IfFontExistsTF{TeX Gyre Heros}{\setsansfont{TeX Gyre Heros}}{\setsansfont{lmsans10-regular.otf}[BoldFont=lmsans10-bold.otf,ItalicFont=lmsans10-oblique.otf]}
\IfFontExistsTF{DejaVu Sans Mono}{\setmonofont{DejaVu Sans Mono}[Scale=MatchLowercase]}{\setmonofont{lmmono10-regular.otf}}
\IfFontExistsTF{Noto Serif CJK SC}{\setCJKmainfont{Noto Serif CJK SC}[AutoFakeSlant=0.15]}{\setCJKmainfont{FandolSong-Regular.otf}[BoldFont=FandolSong-Bold.otf,ItalicFont=FandolKai-Regular.otf]}
\IfFontExistsTF{Noto Sans CJK SC}{\setCJKsansfont{Noto Sans CJK SC}}{\setCJKsansfont{FandolHei-Regular.otf}}
\IfFontExistsTF{Noto Sans Mono CJK SC}{\setCJKmonofont{Noto Sans Mono CJK SC}}{\setCJKmonofont{FandolFang-Regular.otf}}
\usepackage{amsmath,amssymb,amsthm,mathtools,mathrsfs,esint,bm}
\usepackage{booktabs,longtable,array,calc}
\usepackage{enumitem}
\setlist{nosep,leftmargin=1.7em,topsep=3pt}
\usepackage{fvextra}
\DefineVerbatimEnvironment{Highlighting}{Verbatim}{breaklines,breakanywhere,fontsize=\small}
\usepackage{xurl}
\usepackage[unicode,hidelinks,bookmarksnumbered]{hyperref}
\usepackage{bookmark}
\usepackage{fancyhdr}
\pagestyle{fancy}\fancyhf{}
\fancyhead[L]{\small 第一科学主线：近期研究全稿}
\fancyhead[R]{\small 2026-10-07}\fancyfoot[C]{\thepage}
\renewcommand{\headrulewidth}{0.3pt}
\providecommand{\tightlist}{\setlength{\itemsep}{1pt}\setlength{\parskip}{0pt}}
\providecommand{\passthrough}[1]{#1}
\newcommand{\Ric}{\operatorname{Ric}}
\newcommand{\Vol}{\operatorname{Vol}}
\newcommand{\Rm}{\operatorname{Rm}}
\newcommand{\supp}{\operatorname{supp}}
\newcommand{\Var}{\operatorname{Var}}
\newcommand{\tr}{\operatorname{tr}}
\newcommand{\avint}{\fint}
\newcommand{\Jcal}{\mathcal J}
\newcommand{\Ecal}{\mathcal E}
\newcommand{\dd}{\,\mathrm d}
\newcommand{\Rea}{\operatorname{Re}}
\setlength{\parindent}{2em}
\setlength{\parskip}{3pt plus 1pt}
\setlength{\emergencystretch}{3em}
\linespread{1.10}\allowdisplaybreaks[2]
\setlength{\abovedisplayskip}{8pt plus 2pt minus 2pt}
\setlength{\belowdisplayskip}{8pt plus 2pt minus 2pt}
\renewcommand{\arraystretch}{1.12}
\setcounter{tocdepth}{1}\setcounter{secnumdepth}{2}
\pdfstringdefDisableCommands{\def\fint{int}\def\mathcal#1{#1}\def\mathfrak#1{#1}\def\mathscr#1{#1}\def\operatorname#1{#1}\def\varepsilon{epsilon}\def\chi{chi}\def\beta{beta}\def\rho{rho}\def\tau{tau}\def\Phi{Phi}\def\sigma{sigma}\def\kappa{kappa}\def\nabla{D}\def\mathbb#1{#1}\def\({}\def\){}\def\textsuperscript#1{#1}\def\textsubscript#1{#1}}
\title{无密度低幅度估计与全局极小序列接口}
\author{第一科学主线·近期研究分卷}
\date{2026年10月7日}
\begin{document}
\maketitle
\noindent\textbf{原稿编号：}R05。

本章在实际类型投影下证明逆幅度矩、深低幅度能量排除和二次弱测试。随后记录 Anderson 极小序列结果的准确接口与完整端部费用。双重带权引理来自收到的网页补充，不作为本会话独立首创；本章使用它时保留其 \(H^1\) 输入。\cite{R05,Weighted,Anderson}

\section{明确的继承输入}\label{ch5:ux660eux786eux7684ux7ee7ux627fux8f93ux5165}

各流形均闭、光滑，积分用概率测度 \(d\mu_j=dV_j/V_j\)。令 \(E_j\) 为 Hermitian 丛，连接为酉连接，\(L_j=\nabla_j^*\nabla_j\)。

假设已有光滑截面 \(v_j\) 满足 \[
 |v_j|+|L_jv_j|\le C,\qquad
 m_j:=|v_j|^2\longrightarrow m_*>0\quad\text{in }L^2(d\mu_j),
\] \[
 r_j:=L_jv_j-\lambda v_j\longrightarrow0\quad\text{in }L^2(d\mu_j),\qquad \lambda>0.
\tag{5.1}\label{eq:r5:1}
\]

在 O10 中 \(\lambda=2\)，\(m_*=128/81\)。这些是所给有界同丛 resolvent 的前置，并不包含 \(dm_j\to0\)。

进一步，存在区域 \(G_j\)，上面有正交几乎复结构 \(J_j\)，使复化余切丛的正交类型投影 \(\mathsf P_j\) 满足 \[
 N_j:=\mathsf P_j\nabla_jv_j,\qquad
 \int_{G_j}|N_j|^2\,d\mu_j\longrightarrow0,
\tag{5.2}\label{eq:r5:2}
\] 以及 \[
 \int_{M_j\setminus G_j}|\nabla_jv_j|^2\,d\mu_j\longrightarrow0.
\tag{5.3}\label{eq:r5:3}
\]

这里 \(\mathsf P=(I-iJ)/2\)；换用共轭类型不改变本证明。

对所给 O10 输入，可取 \(G_j=\{\rho_j>c\}\)，其中 \(0<c<\rho_*=4/3\) 固定。\eqref{eq:r5:2} 是已有首阶投影。\eqref{eq:r5:3} 不需要 \(L^4\) 梯度估计：以 \(\chi(\rho_j)^2v_j\) 测试 \(L_jv_j\)，得 \[
 \int\chi^2|\nabla v_j|^2\,d\mu_j
 \le C\int(\chi^2+|d\chi|^2)\,d\mu_j=o(1),
\] 其中 \(\chi=1\) 于 \(\rho\le c\)，在另一个固定的、低于 \(\rho_*\) 的阈值之前变为0。右端使用已有 \(\rho_j\to\rho_*\) 和 \(d\rho_j\to0\)。此步不交换协变导数。

记 \(q_j=|\nabla_jv_j|^2\)。由 \eqref{eq:r5:1} 配对 \(v_j\) 可知 \[
 \int q_j\,d\mu_j\longrightarrow\lambda m_*.
\tag{5.4}\label{eq:r5:4}
\]

\section{逆幅度矩估计}\label{ch5:ux9006ux5e45ux5ea6ux77e9ux4f30ux8ba1}

\textbf{命题。} 在 \eqref{eq:r5:1}---\eqref{eq:r5:3} 下，对每个固定 \(0<p<1\)、\(\varepsilon>0\)， \[
 \boxed{
 \limsup_{j\to\infty}\int
 \frac{q_j}{(m_j+\varepsilon)^p}\,d\mu_j
 \le\frac{\lambda m_*}{(1-p)(m_*+\varepsilon)^p}.
 }
\tag{5.5}\label{eq:r5:5}
\]

常数不使用曲率、直径、非塌缩、Sobolev 常数或 Hessian 预算。顺序是：先固定 \(p,\varepsilon\)，再令 \(j\to\infty\)。不声称对任意事先指定的 \(\varepsilon_j\to0\) 有同一有限误差估计。

\subsection{1 类型投影的点态代数}\label{ch5:ux7c7bux578bux6295ux5f71ux7684ux70b9ux6001ux4ee3ux6570}

在 \(G_j\) 上写 \(\nabla v=N+A\)，其中 \(N=\mathsf P\nabla v\)，\(A=(I-\mathsf P)\nabla v\)。二者正交，\(q=|N|^2+|A|^2\)。

一个纯 \((1,0)\) 或 \((0,1)\) 复一形式 \(\alpha\) 满足 \(|2\operatorname{Re}\alpha|^2=2|\alpha|^2\)。因此对任何 \(\delta>0\)， \[
 |dm|^2
 \le 2(1+\delta)m q+C_\delta m|N|^2.
\tag{5.6}\label{eq:r5:6}
\] 这是对 \(dm=2\operatorname{Re}\langle v,\nabla v\rangle\) 使用加权 Young，不是导数估计。在补集上仍有一般 Kato 界 \[
 |dm|^2\le4mq.
\tag{5.7}\label{eq:r5:7}
\]

\subsection{2 在原闭流形上测试}\label{ch5:ux5728ux539fux95edux6d41ux5f62ux4e0aux6d4bux8bd5}

用光滑截面 \((m+\varepsilon)^{-p}v\) 测试 \(Lv\)。精确地， \[
 \int\frac q{(m+\varepsilon)^p}
 =\operatorname{Re}\int\frac{\langle Lv,v\rangle}{(m+\varepsilon)^p}
 +\frac p2\int\frac{|dm|^2}{(m+\varepsilon)^{p+1}}.
\tag{5.8}\label{eq:r5:8}
\] 所有积分仍为归一化积分。

将第一项写为 \(F_j\)，并把左端分成 \(I_{G,j}+I_{G^c,j}\)。由 \eqref{eq:r5:6}---\eqref{eq:r5:7}，只要 \(p(1+\delta)<1\)， \[
 [1-p(1+\delta)]I_{G,j}
 \le F_j+|2p-1|I_{G^c,j}
   +C_{p,\delta}\varepsilon^{-p}\int_{G_j}|N_j|^2.
\tag{5.9}\label{eq:r5:9}
\]

固定 \(\varepsilon\) 后，\eqref{eq:r5:3} 使 \(I_{G^c,j}\to0\)，\eqref{eq:r5:2} 使最后一项趋0。由有界幅度、\eqref{eq:r5:1} 和连续函数收敛， \[
 F_j\longrightarrow\frac{\lambda m_*}{(m_*+\varepsilon)^p}.
\] 故 \[
 \limsup_j I_{G,j}\le
 \frac{\lambda m_*}
 {[1-p(1+\delta)](m_*+\varepsilon)^p}.
\] 最后令 \(\delta\downarrow0\)，得到 \eqref{eq:r5:5}。没有在有限误差中任意指定 \(\delta_j\) 或 \(\varepsilon_j\)。

\section{深低幅度区的能量排除}\label{ch5:ux6df1ux4f4eux5e45ux5ea6ux533aux7684ux80fdux91cfux6392ux9664}

对每个固定 \(0<t<m_*\)，由 \eqref{eq:r5:5}， \[
 \begin{aligned}
 \limsup_j\int_{\{m_j\le t\}}q_j
 &\le (t+\varepsilon)^p
       \limsup_j\int\frac{q_j}{(m_j+\varepsilon)^p}\\
 &\le\frac{\lambda m_*}{1-p}
       \left(\frac{t+\varepsilon}{m_*+\varepsilon}\right)^p.
 \end{aligned}
\] 令 \(\varepsilon\downarrow0\)，得到 \[
 \boxed{
 \limsup_j\int_{\{m_j\le t\}}q_j\,d\mu_j
 \le\frac{\lambda m_*}{1-p}\left(\frac t{m_*}\right)^p,
 \qquad 0<p<1.
 }
\tag{5.10}\label{eq:r5:10}
\]

当 \(0<t<m_*/e\) 时，取固定的 \[
 p=1-\frac1{\log(m_*/t)},
\] 便有 \[
 \boxed{
 \limsup_j\int_{\{m_j\le t\}}q_j\,d\mu_j
 \le e\lambda\,t\log\frac{m_*}{t}.
 }
\tag{5.11}\label{eq:r5:11}
\]

在 O10 中 \(\lambda=2\)。特别地， \[
 \lim_{t\downarrow0}\limsup_j
 \int_{\{m_j\le t\}}|\nabla v_j|^2\,d\mu_j=0.
\tag{5.12}\label{eq:r5:12}
\]

\subsection{1 不能从这里推出什么}\label{ch5:ux4e0dux80fdux4eceux8fd9ux91ccux63a8ux51faux4ec0ux4e48}

\eqref{eq:r5:12} 只控制模长趋于零的最深井。它不证明对固定 \(0<t<m_*\)，\(\int_{m_j\le t}q_j\to0\)；不证明 \(dm_j\to0\) 于 \(L^2\)；不证明 \(q_j\) 在所有小体积集合上一致可积；也不证明第五线丛截面的非零质量。

固定中间幅度层 \(a<m_j<b<m_*\) 的体积可以趋零，但这里仍可能有非零径向能量。将这个区别省去，就会再次把一个合法中间估计误写成 O10 已解。

\subsection{2 为何这类矩估计不能单独消灭全部径向缺陷}\label{ch5:ux4e3aux4f55ux8fd9ux7c7bux77e9ux4f30ux8ba1ux4e0dux80fdux5355ux72ecux6d88ux706dux5168ux90e8ux5f84ux5411ux7f3aux9677}

取振幅变量 \(t\in[0,m_*]\)，考虑纯形式的能量测度 \[
 d\nu_q(t)=\lambda\,dt,\qquad d\nu_R(t)=2\lambda t\,dt.
\tag{5.13}\label{eq:r5:13}
\] 对光滑测试 \(F\)，它们满足 \[
 \int F\,d\nu_q+\frac12\int F'\,d\nu_R
 =\lambda m_* F(m_*),
\] 与 \eqref{eq:r5:8} 的极限质量关系相容，并在 \(\varepsilon\downarrow0\) 时达到 \eqref{eq:r5:5} 的逆幅度矩常数。

O10 取 \(\lambda=2\) 时，\(d\nu_R=4t\,dt\) 正是前轮文件已显示的径向预算。它还满足 \eqref{eq:r5:12}，但总径向能量 \(\nu_R([0,m_*])=\lambda m_*^2\) 非零。

这不是原方程的实际反例。它只严格说明：已用的类型代数、标量质量关系及其逆幅度矩，并不足以排除全部径向集中。不能再把同类标量权重的形式变化当作自动能够闭合的证明路线。

\section{一个可以直接接入弱测试的二次能量关系}\label{ch5:ux4e00ux4e2aux53efux4ee5ux76f4ux63a5ux63a5ux5165ux5f31ux6d4bux8bd5ux7684ux4e8cux6b21ux80fdux91cfux5173ux7cfb}

若线丛测试 \(w_j\) 满足 \[
 \|w_j\|_\infty+\|L_ww_j\|_\infty
 +\|\nabla w_j\|_{4,\mathrm{av}}\le C,
\] 则 \[
 \boxed{
 \int(q_j-\lambda m_*)|w_j|^2\,d\mu_j\longrightarrow0.
 }
\tag{5.14}\label{eq:r5:14}
\]

证明：由 \[
 \Delta m_j=2\lambda m_j-2q_j
       +2\operatorname{Re}\langle r_j,v_j\rangle
\] 及 Stokes，左边等于 \[
 \lambda\int(m_j-m_*)|w_j|^2
 +\int\operatorname{Re}\langle r_j,v_j\rangle|w_j|^2
 -\frac12\int(m_j-m_*)\Delta|w_j|^2.
\] 其中 \[
 \Delta|w_j|^2
 =2\operatorname{Re}\langle L_ww_j,w_j\rangle
       -2|\nabla w_j|^2
\] 有一致 \(L^2\) 界，三项均趋0。

\eqref{eq:r5:14} 不需要 \(q_j\) 的一致 \(L^2\) 或全局一致可积性，但也不把 \(q_j\) 等同于第五线丛的 \(|c_j|^2\)。两者之间的径向项仍须保留。

\section{O11：Anderson 接口的核对与仍缺的下界}\label{ch5:o11anderson-ux63a5ux53e3ux7684ux6838ux5bf9ux4e0eux4ecdux7f3aux7684ux4e0bux754c}

外部核对来源：M. T. Anderson, \emph{Canonical Metrics on 3-Manifolds and 4-Manifolds}, arXiv:math/0511034，Theorem 4.10，(4.43)、(4.45)，以及 (4.65) 后的讨论。

论文提供特殊选择的理想极小序列：紧轨形情形有 ALE 项参与的能量等式；开放/塌缩情形有相应下半连续不等式。它没有证明任意给定近极小序列自动具有该结构，也没有证明原异常类和连接柄全部进入指定的完整 AE 核心。

其 (4.65) 明确允许沿球面空间形式、而不只是沿三维球面作广义连通和；随后专门指出开放极限域与原流形拓扑的关系仍需研究。本文引用这一明确范围，不以 2005 年的讨论冒充 2026 年前沿的穷尽结论。

\subsection{1 核心费用需要保留端部修正}\label{ch5:ux6838ux5fc3ux8d39ux7528ux9700ux8981ux4fddux7559ux7aefux90e8ux4feeux6b63}

对满足相应渐近条件、有限拓扑型、负定交叉形式的完整 ALE 四流形 \(N\)，记 \(e\) 为端数，\(r=b_1(N)\)，\(m=b_2(N)\)，\(G=\sum_a|\Gamma_a|^{-1}\)。以 \[
 \tau(N)=\frac1{12\pi^2}\int(W_+^2-W_-^2)-\eta_\infty
\] 固定 \(\eta_\infty\) 的符号。

若无其他边界项残留，相对 Gauss--Bonnet、signature 及 \(\chi(N)=2-2r+m-e\) 给出精确式 \[
 \begin{aligned}
 P_N={}&4\int|W^+|^2+\frac1{12}\int s^2
 +8\pi^2m+32\pi^2r\\
 &+16\pi^2(e-2+G)-24\pi^2\eta_\infty.
 \end{aligned}
\tag{5.15}\label{eq:r5:15}
\]

AE 情形 \(G=e,\eta_\infty=0\) 才给出 \[
 P_N\ge8\pi^2m+32\pi^2(r+e-1)+\frac1{12}\int s^2.
\tag{5.16}\label{eq:r5:16}
\]

这些费用与用户新稿一致，本文不将它们重新登记为新结果。其意义是：拥有一个二周期，不自动拥有 \(8\pi^2\) 的无迹 Ricci 成本。Eguchi--Hanson 的 \(m=1,e=1,G=1/2,\eta_\infty=0\) 使 \eqref{eq:r5:15} 的拓扑部分为0，符合它的 Ricci-flat 性。

\subsection{极小序列的标记与费用}\label{ch5:ux6781ux5c0fux5e8fux5217ux7684ux6807ux8bb0ux4e0eux8d39ux7528}

Anderson 结果提供的是可选择极小序列的结构和相应能量方向；AE 费用需要完整欧氏端及所列拓扑输入，一般 ALE 费用需要保留端群和谱修正。后来的严格次四构造已经改变原族的候选最优值，因此不把``所有手术标记都按标准费用独立收费''作为已知结论。\cite{System}

\section{双重带权输入及其与本章的连接}\label{ch5:ux53ccux91cdux5e26ux6743ux8f93ux5165ux53caux5176ux4e0eux672cux7ae0ux7684ux8fdeux63a5}

设 \(v\) 满足 \(|v|+|\nabla^*\nabla v|\le C\)。对光滑张量 \(e\)，以 \(|e|^2v\) 测试，得到 \[
\int|e|^2|\nabla v|^2
=\operatorname{Re}\int|e|^2\langle\nabla^*\nabla v,v\rangle
-2\operatorname{Re}\int |e|\langle d|e|\otimes v,\nabla v\rangle.
\] Kato 与 Young 不等式给 \[
\||e|\nabla v\|_2\le C\|e\|_{H^1}.
\] 若线丛测试 \(w\) 具有有界幅度、粗 Laplacian 和二阶 \(L^2\) 预算，取 \(e=\nabla w\) 即得 \[
\||\nabla v||\nabla w|\|_2\le C.
\] 对全局有界收缩张量 \(\mathcal C\)，若 \(|\nabla\mathcal C|\le C(1+|\nabla\Phi|+|\nabla v|)\)，且 \(\zeta\to0\) 于 \(H^1\)，一次分部积分使 \[
\left|\fint\zeta\,\mathcal C\,\nabla^2v\,\nabla w\right|
\le C\|\zeta\|_{H^1}\to0.
\] 这是外部网页补充的输入。其系数需要 \(H^1\) 小性；只知 \(m-m_*\to0\) 于 \(L^2\) 时，不可把 \(\zeta\) 直接换成 \(m-m_*\)。本章的逆幅度估计和二次弱测试补充其适用范围，但不将中间幅度层的径向能量宣称为零。\cite{Weighted}

\begin{thebibliography}{P99}
\small
\bibitem[R05]{R05} 本会话研究原稿，\emph{无密度低幅度估计与全局极小序列接口}，2026-09-30。\path{O10_Low_Amplitude_Energy_and_O11_Scope.md}。

\bibitem[P06]{Weighted} 所提供网页续研 E，\emph{A double weighted test removes a class of uncontrolled negative Hessians}，2026-09-30。

\bibitem[E09]{Anderson} M. T. Anderson, \emph{Canonical Metrics on 3-Manifolds and 4-Manifolds}, arXiv:math/0511034, Theorem 4.10, (4.43), (4.45). \url{https://arxiv.org/abs/math/0511034}.

\bibitem[P08]{System} S1003A01、S1003A02，\emph{Mathematical Synthesis}，2026-10-07；\path{20261007_Einstein_Gaps_Two_Sessions(1).zip} 中 \path{01_Mathematical_Synthesis.md}。

\end{thebibliography}
\end{document}
